\documentclass[11pt]{article}
\usepackage[margin=1in]{geometry}
\usepackage[T1]{fontenc}
\usepackage[utf8]{inputenc}
\usepackage{lmodern}
\usepackage{microtype}
\usepackage[table]{xcolor}
\usepackage{amsmath,amssymb,amsfonts,mathtools,bm}
\IfFileExists{physics.sty}{%
  \usepackage{physics}%
}{%
  % Portable subset used by this project when the optional physics package
  % is not installed in the local TeX distribution.
  \NewDocumentCommand{\pdv}{o m m}{%
    \IfNoValueTF{##1}{\frac{\partial ##2}{\partial ##3}}%
    {\frac{\partial^{##1} ##2}{\partial ##3^{##1}}}%
  }
  \NewDocumentCommand{\dv}{o m m}{%
    \IfNoValueTF{##1}{\frac{\mathrm{d} ##2}{\mathrm{d} ##3}}%
    {\frac{\mathrm{d}^{##1} ##2}{\mathrm{d} ##3^{##1}}}%
  }
  \providecommand{\dd}{\mathop{}\!\mathrm{d}}
  \providecommand{\grad}{\nabla}
  \providecommand{\abs}[1]{\left\lvert ##1\right\rvert}
  \providecommand{\norm}[1]{\left\lVert ##1\right\rVert}
  \providecommand{\ket}[1]{\left\lvert ##1\right\rangle}
  \providecommand{\bra}[1]{\left\langle ##1\right\rvert}
}
\usepackage{booktabs,array,longtable,tabularx}
\usepackage{hyperref}
\usepackage{enumitem}
\usepackage{fancyhdr}
\usepackage{titlesec}
\usepackage{tcolorbox}
\usepackage{ragged2e}
\usepackage{setspace}
\IfFileExists{siunitx.sty}{%
  \usepackage{siunitx}%
}{%
  % Minimal text-safe fallback for the small number of \SI and \si calls in
  % this repository. Full siunitx formatting is used when available.
  \providecommand{\si}[1]{\ensuremath{\mathrm{##1}}}
  \providecommand{\SI}[2]{\ensuremath{##1\,\mathrm{##2}}}
}
\hypersetup{colorlinks=true, linkcolor=blue!50!black, urlcolor=blue!50!black, citecolor=blue!50!black}
\definecolor{TitleBlue}{HTML}{163A5F}
\definecolor{AccentTeal}{HTML}{2D6F73}
\definecolor{SoftBlue}{HTML}{EAF3F8}
\definecolor{SoftTeal}{HTML}{EDF8F6}
\definecolor{SoftGray}{HTML}{F5F7FA}
\definecolor{RuleGray}{HTML}{D7DEE8}
\pagestyle{fancy}
\fancyhf{}
\lhead{RadiationEffects\_proj2}
\rhead{Technical Notes}
\cfoot{\thepage}
\setlength{\headheight}{14pt}
\setlength{\parskip}{0.5em}
\setlength{\parindent}{0pt}
\onehalfspacing
\numberwithin{equation}{section}
\titleformat{\section}{\Large\bfseries\color{TitleBlue}}{\thesection}{0.6em}{}
\titleformat{\subsection}{\large\bfseries\color{AccentTeal}}{\thesubsection}{0.6em}{}
\titleformat{\subsubsection}{\normalsize\bfseries\color{TitleBlue}}{\thesubsubsection}{0.6em}{}
\newtcolorbox{overviewbox}{colback=SoftBlue,colframe=TitleBlue,boxrule=0.8pt,arc=2mm,left=3mm,right=3mm,top=2mm,bottom=2mm}
\newtcolorbox{physicsbox}{colback=SoftTeal,colframe=AccentTeal,boxrule=0.8pt,arc=2mm,left=3mm,right=3mm,top=2mm,bottom=2mm}
\newtcolorbox{notebox}{colback=SoftGray,colframe=AccentTeal,boxrule=0.6pt,arc=2mm,left=3mm,right=3mm,top=2mm,bottom=2mm}
\newcommand{\DocTitle}[3]{%
\begin{center}
{\LARGE \bfseries \color{TitleBlue} #1}\\[0.5em]
{\large \color{AccentTeal} #2}\\[0.8em]
{\normalsize #3}
\end{center}
\vspace{0.5em}}
\newenvironment{symtable}{\rowcolors{2}{SoftGray}{white}\begin{longtable}{>{\RaggedRight\arraybackslash}p{0.18\textwidth} >{\RaggedRight\arraybackslash}p{0.25\textwidth} >{\RaggedRight\arraybackslash}p{0.47\textwidth}}\toprule \textbf{Symbol / acronym} & \textbf{Name} & \textbf{Meaning in this document} \\ \midrule\endfirsthead\toprule \textbf{Symbol / acronym} & \textbf{Name} & \textbf{Meaning in this document} \\ \midrule\endhead}{\bottomrule\end{longtable}}


\newcommand{\ProjectTOC}{%
\vspace{0.6em}
{\hypersetup{linkcolor=TitleBlue,urlcolor=TitleBlue,citecolor=TitleBlue}\tableofcontents}
\vspace{1.0em}}

\newcommand{\SymbolsAcronymsSection}{%
\section*{Symbols and acronyms used in this document}%
\addcontentsline{toc}{section}{Symbols and acronyms used in this document}}

\newcommand{\rvec}{\bm{r}}
\newcommand{\qvec}{\bm{q}}
\newcommand{\nvec}{\bm{n}}
\newcommand{\xvec}{\bm{x}}
\newcommand{\kB}{k_{\mathrm{B}}}
\newcommand{\qp}{\mathrm{qp}}
\newcommand{\JJ}{\mathrm{JJ}}
\newcommand{\mfp}{\Lambda}
\allowdisplaybreaks

\newcommand{\kB}{k_{\mathrm{B}}}
\newcommand{\vvg}{\mathbf{v}_g}
\newcommand{\qvec}{\mathbf{q}}
\newcommand{\rvec}{\mathbf{r}}
\newcommand{\absval}[1]{\left\lvert #1 \right\rvert}
\rhead{Module 4b Physics Note}

\begin{document}

\DocTitle{Module 4b: Two-Dimensional Phonon-Aware Thermal Transport}{Derivation-Focused Physics Note for \texttt{RadiationEffects\_proj2}}{Thermal transport model with defect-sensitive phonon-limited conductivity}

\hrule
\vspace{0.8em}

\begin{overviewbox}
\textbf{Purpose of this note.} Module 4a treated heat flow in silicon with the continuum heat equation and a prescribed thermal conductivity. Module 4b keeps the same macroscopic temperature field description, but derives a \emph{phonon-aware} constitutive law in which the thermal conductivity changes because radiation-induced defects scatter the phonons that carry heat. The purpose is not to jump immediately to a full phonon Boltzmann solver. The purpose is to build the physically justified intermediate model that sits between bulk Fourier conduction and a full microscopic transport theory.
\end{overviewbox}

\ProjectTOC


\SymbolsAcronymsSection
\renewcommand{\arraystretch}{1.2}
\begin{longtable}{p{0.17\textwidth} p{0.26\textwidth} p{0.49\textwidth}}
\toprule
\textbf{Symbol} & \textbf{Name} & \textbf{Meaning} \\
\midrule
\endfirsthead
\toprule
\textbf{Symbol} & \textbf{Name} & \textbf{Meaning} \\
\midrule
\endhead
$T(\rvec,t)$ & temperature field & Coarse-grained thermal state variable used in the reduced continuum description. \\
$\rho$ & mass density & Mass per unit volume of the silicon medium. \\
$c_p$ & specific heat & Heat capacity per unit mass at constant pressure. \\
$u$ & internal energy density & Thermal energy stored per unit volume. \\
$Q(\rvec,t)$ & volumetric heat source & Heat generated per unit volume per unit time. \\
$\qvec$ & heat flux vector & Thermal energy flow per unit area per unit time. \\
$k$ & thermal conductivity & Constitutive coefficient relating heat flux to temperature gradient. In Module 4b, $k=k(C,T)$. \\
$C(\rvec,t)$ & defect concentration field & Damage variable supplied by Module 3 and used to modify phonon scattering and conductivity. \\
$f$ & phonon distribution function & Occupation of phonon states in position, direction, frequency, branch, and time. \\
$f^{(0)}$ & equilibrium phonon distribution & Local Bose-Einstein equilibrium distribution. \\
$\delta f$ & nonequilibrium correction & Small departure from local equilibrium induced by gradients. \\
$\vvg$ & phonon group velocity & Velocity at which phonon wave packets transport energy. \\
$\tau$ & relaxation time & Characteristic time for scattering to relax the phonon distribution toward equilibrium. \\
$\ell$ & mean free path & Average distance a phonon travels between significant scattering events. \\
$C_v$ & volumetric heat capacity & Heat capacity per unit volume used in kinetic-theory conductivity estimates. \\
$\tau_{\mathrm{pp}}$ & phonon-phonon scattering time & Effective scattering time associated with anharmonic phonon interactions. \\
$\tau_{\mathrm{b}}$ & boundary scattering time & Effective scattering time associated with boundaries or interfaces. \\
$\tau_{\mathrm{d}}$ & defect scattering time & Effective scattering time associated with radiation-induced defects. \\
$\tau_{\mathrm{eff}}$ & effective relaxation time & Combined relaxation time resulting from multiple scattering channels. \\
$k_{\mathrm{bulk}}(T)$ & bulk conductivity & Thermal conductivity of undamaged silicon as a function of temperature. \\
$\Psi(C,T)$ & conductivity reduction factor & Dimensionless factor used to reduce bulk conductivity because of defects. \\
$\alpha(T),\beta(T)$ & reduction parameters & Model parameters controlling how strongly defects suppress conductivity. \\
$D(T)$ & defect diffusion coefficient & Temperature-dependent diffusion coefficient used in Module 3. \\
$D_0$ & diffusion prefactor & Pre-exponential factor in the Arrhenius diffusion law. \\
$E_a$ & activation energy & Energy barrier for defect diffusion. \\
$k_{\mathrm{ann}}(T)$ & annealing rate & Temperature-dependent reaction or annealing rate. \\
$k_0$ & annealing prefactor & Pre-exponential factor for annealing kinetics. \\
$E_{\mathrm{ann}}$ & annealing activation energy & Energy barrier for annealing or reaction kinetics. \\
$\kB$ & Boltzmann constant & Constant relating temperature to thermal energy scale. \\
$L$ & characteristic length & Representative device length scale used in time-scale estimates. \\
$\alpha_{\mathrm{th},0}$ & nominal thermal diffusivity & Reference thermal diffusivity defined as $k_0/(\rho c_p)$ for scaling analysis. \\
\bottomrule
\end{longtable}

\section{Physical role of Module 4b in the full project}

The full project aims to connect a radiation event in silicon to long-term device degradation. In that chain, the thermal module matters because temperature changes defect mobility, reaction rates, annealing rates, and eventually electrical behavior. Module 4a already captured the statement ``temperature matters.'' Module 4b goes further and captures the statement ``the defect population itself can modify heat transport.''

That extra statement is physically important in silicon because heat is carried primarily by lattice vibrations, that is, by phonons. Radiation-induced defects disturb the crystal lattice, and that disturbance increases phonon scattering. Increased phonon scattering shortens the average distance a phonon travels before losing memory of its previous state. Since thermal conductivity depends strongly on that transport distance, the thermal conductivity can degrade as defect concentration increases.

Thus the feedback chain for Module 4b is
\begin{equation}
\text{radiation} \rightarrow \text{defects} \rightarrow \text{phonon scattering} \rightarrow k \text{ changes} \rightarrow T(\rvec,t) \text{ changes} \rightarrow \text{defect evolution changes.}
\end{equation}

This feedback does not exist in Module 4a if $k$ is held fixed. That is the conceptual reason for splitting Module 4 into Module 4a and Module 4b.

\section{Why the plain continuum model can become insufficient}

The continuum heat equation used in Module 4a is
\begin{equation}
\rho c_p \frac{\partial T}{\partial t} = \nabla\cdot\left(k\nabla T\right) + Q,
\end{equation}
where $\rho$ is the mass density, $c_p$ is the specific heat at constant pressure, $T$ is temperature, $k$ is thermal conductivity, and $Q$ is the volumetric heat source.

This equation is very useful, but it silently assumes several pieces of physics:
\begin{enumerate}[label=\arabic*.]
    \item At each point in space, a local temperature is well defined.
    \item Heat transport is diffusive rather than ballistic.
    \item The microscopic carriers equilibrate rapidly enough that a constitutive law $\qvec=-k\nabla T$ is valid.
    \item The conductivity $k$ can be treated as either constant or as a simple prescribed function independent of the evolving radiation damage state.
\end{enumerate}

For many engineering-scale problems, those assumptions are adequate. But the present project is not only a thermal problem; it is a radiation-damage problem. Since the damage is one of the evolving state variables, it is physically natural to ask whether the same damage that changes electrical behavior also changes thermal transport.

That question points directly to phonon physics.

\section{Heat transport in crystalline silicon as phonon transport}

\subsection{Why phonons appear}

A silicon crystal consists of atoms arranged in a periodic lattice. The atoms are not fixed absolutely in space. They vibrate around equilibrium positions. Because the atoms are coupled to their neighbors, those vibrations propagate through the crystal as collective normal modes. In quantum language, the quanta of those lattice vibrational modes are called \textbf{phonons}.

One should not think of a phonon as a tiny classical particle in the naive sense. A phonon is a quantized normal mode of lattice vibration. However, for transport theory, it is extremely useful to treat phonons as quasiparticles carrying energy, momentum, and a group velocity. In that picture, heat flow in a non-metallic crystal emerges from the transport of a population of phonons moving and scattering throughout the material.

Silicon is not a simple metal in which free electrons dominate thermal conductivity under ordinary device conditions. In crystalline silicon, lattice thermal conduction is dominant over electronic thermal conduction in many regimes of interest. Therefore, a damage-induced change in phonon transport can produce a meaningful change in the overall thermal conductivity.

\subsection{Local equilibrium versus distribution-function language}

In Module 4a, the thermal state was represented by a single scalar field $T(\rvec,t)$. In a more microscopic description, one instead introduces a phonon distribution function,
\begin{equation}
 f = f(\rvec,\mathbf{s},\omega,p,t),
\end{equation}
where
\begin{itemize}
    \item $\rvec$ is position,
    \item $\mathbf{s}$ is propagation direction,
    \item $\omega$ is angular frequency,
    \item $p$ labels the phonon branch or polarization,
    \item $t$ is time.
\end{itemize}

The quantity $f$ tells us how strongly each phonon state is occupied. If the system is in local thermal equilibrium, $f$ reduces locally to the Bose-Einstein equilibrium distribution associated with the local temperature. If the system is not in local equilibrium, then a single temperature field is no longer enough to describe the full thermal state.

Module 4b does \emph{not} attempt to solve for the full distribution function. Instead, it asks: what reduced model can we derive from phonon transport ideas that is still practical for the larger multiphysics project?

\section{Microscopic starting point: the phonon Boltzmann transport equation}

A standard microscopic starting point is the phonon Boltzmann transport equation (BTE):
\begin{equation}
\frac{\partial f}{\partial t} + \vvg\cdot \nabla f = \left(\frac{\partial f}{\partial t}\right)_{\!\mathrm{coll}},
\end{equation}
where $\vvg$ is the phonon group velocity and the right-hand side represents collisions or scattering.

This equation says the following. The left-hand side describes free transport of the distribution function through phase space. The right-hand side describes how scattering drives the phonon population away from or toward local equilibrium.

A common simplification is the relaxation-time approximation:
\begin{equation}
\left(\frac{\partial f}{\partial t}\right)_{\!\mathrm{coll}} = -\frac{f-f^{(0)}(T)}{\tau},
\end{equation}
where $f^{(0)}(T)$ is the local equilibrium distribution and $\tau$ is an effective relaxation time. This approximation says that scattering tends to relax the actual distribution toward the local equilibrium distribution over a characteristic time scale $\tau$.

The usefulness of this equation for the present project is not that we plan to solve it directly in the first implementation. Its usefulness is that it shows \emph{where the thermal conductivity comes from}. Thermal conductivity is not just an arbitrary constant. It is a macroscopic summary of how the phonon distribution responds to a temperature gradient under transport and scattering.

\section{From the BTE to Fourier-like heat conduction}

\subsection{Linear-response picture}

Suppose the crystal is close to local equilibrium, with only weak spatial temperature gradients. Then the actual phonon distribution can be written as
\begin{equation}
 f = f^{(0)} + \delta f,
\end{equation}
where $\delta f$ is a small correction. Substituting this into the relaxation-time BTE and keeping only first-order corrections gives, schematically,
\begin{equation}
 \delta f \propto -\tau\, \vvg\cdot \nabla f^{(0)}.
\end{equation}

Because $f^{(0)}$ depends on temperature, and temperature depends on position, we have by the chain rule
\begin{equation}
\nabla f^{(0)} = \frac{\partial f^{(0)}}{\partial T}\nabla T.
\end{equation}
Therefore,
\begin{equation}
 \delta f \propto -\tau\, (\vvg\cdot \nabla T)\frac{\partial f^{(0)}}{\partial T}.
\end{equation}

This result already contains the basic physics: a temperature gradient slightly biases the phonon population away from isotropic equilibrium, and the amount of bias is controlled by the time available before scattering randomizes the phonon state.

\subsection{Heat flux from the perturbed distribution}

The heat flux is obtained by summing the energy carried by each phonon state weighted by its transport velocity. Symbolically,
\begin{equation}
 \qvec = \sum_{p}\int \hbar\omega\, \vvg\, \delta f\, g(\omega,p)\, d\omega\, d\Omega,
\end{equation}
where $g(\omega,p)$ is a density-of-states factor and $d\Omega$ represents integration over propagation directions.

Inserting the linearized correction $\delta f$ produces a constitutive law of the form
\begin{equation}
 \qvec = -k\nabla T.
\end{equation}
Thus Fourier's law is not separate from phonon transport. It is the near-equilibrium, diffusive-limit consequence of phonon transport.

\subsection{Kinetic-theory estimate for the conductivity}

A compact and physically revealing estimate is
\begin{equation}
 k \sim \frac{1}{3} C_v v \ell,
\end{equation}
where
\begin{itemize}
    \item $C_v$ is the volumetric heat capacity associated with the phonons,
    \item $v$ is a characteristic phonon group speed,
    \item $\ell$ is an effective phonon mean free path.
\end{itemize}

This relation is not exact in all circumstances, but it captures the most important project-level message: \textbf{thermal conductivity is proportional to the phonon mean free path}. If radiation damage reduces $\ell$, then the conductivity decreases.

Since mean free path and relaxation time are related by
\begin{equation}
\ell = v\tau,
\end{equation}
we may also write, schematically,
\begin{equation}
 k \sim \frac{1}{3} C_v v^2 \tau.
\end{equation}
This makes the role of scattering even more explicit. Stronger scattering means smaller $\tau$, hence smaller $k$.

\section{Scattering physics relevant to this project}

\subsection{Phonon-phonon scattering}

Even in a perfect crystal, phonons scatter from one another through anharmonic interactions. This scattering becomes stronger as temperature increases because the phonon population increases. Phonon-phonon scattering is one of the reasons thermal conductivity is temperature dependent even in a nominally defect-free crystal.

\subsection{Phonon-boundary scattering}

If the device dimensions become comparable to important phonon mean free paths, then boundaries and interfaces matter. Rough boundaries can randomize phonon propagation directions, reducing the effective transport length and therefore reducing thermal conductivity relative to the bulk value.

\subsection{Phonon-defect scattering}

This is the key mechanism for Module 4b. Radiation-induced defects disturb the ideal periodicity of the silicon lattice. That disturbance can arise through vacancies, interstitials, clusters, complexes, or strain fields around defect sites. Once the lattice is no longer perfectly periodic, phonons scatter from those inhomogeneities.

The exact microscopic scattering rate depends on defect type, size, concentration, and the phonon wavelength. For the project-level reduced model, we do not need the full spectral scattering law at the outset. What we do need is the qualitative and semi-quantitative fact that more defects generally imply more scattering, and more scattering generally implies lower thermal conductivity.

\section{Combining scattering channels: Matthiessen-like reasoning}

A common reduced transport idea is that scattering rates add:
\begin{equation}
\frac{1}{\tau_{\mathrm{eff}}} = \frac{1}{\tau_{\mathrm{pp}}} + \frac{1}{\tau_{\mathrm{b}}} + \frac{1}{\tau_{\mathrm{d}}} + \cdots,
\end{equation}
where
\begin{itemize}
    \item $\tau_{\mathrm{pp}}$ is the phonon-phonon scattering time,
    \item $\tau_{\mathrm{b}}$ is the boundary scattering time,
    \item $\tau_{\mathrm{d}}$ is the defect scattering time.
\end{itemize}

This is a Matthiessen-type approximation. It is not exact for all spectral details, but it is a powerful reduced model because it translates new scattering channels into a new effective relaxation time.

If the defects become more numerous, then $\tau_{\mathrm{d}}$ decreases, so $1/\tau_{\mathrm{d}}$ increases. Therefore $\tau_{\mathrm{eff}}$ decreases. Since $k\propto \tau_{\mathrm{eff}}$, the conductivity decreases.

\subsection{A simple defect-dependent conductivity model}

One practical project-level model is
\begin{equation}
 k(C,T) = k_{\mathrm{bulk}}(T)\, \Psi(C,T),
\end{equation}
where
\begin{itemize}
    \item $C$ is a suitable defect concentration field from Module 3,
    \item $k_{\mathrm{bulk}}(T)$ is the conductivity of undamaged silicon,
    \item $\Psi(C,T)$ is a reduction factor satisfying $0<\Psi\le 1$.
\end{itemize}

A simple first model may be chosen as
\begin{equation}
 \Psi(C,T) = \frac{1}{1+\alpha(T) C},
\end{equation}
or, in exponential form,
\begin{equation}
 \Psi(C,T) = e^{-\beta(T) C}.
\end{equation}

These are not fundamental microscopic laws. They are reduced constitutive closures motivated by the fact that increasing defect concentration should increase phonon scattering and reduce conductivity. Their parameters $\alpha(T)$ or $\beta(T)$ would later be calibrated against literature, molecular dynamics, experiment, or higher-fidelity transport calculations.

\section{Derivation of the Module 4b governing equation}

\subsection{Energy conservation in a solid}

We begin, as in Module 4a, with local energy conservation. Let $u(\rvec,t)$ be the thermal internal energy density, and let $Q(\rvec,t)$ be a volumetric heat source. Then
\begin{equation}
 \frac{\partial u}{\partial t} + \nabla\cdot \qvec = Q.
\end{equation}

If local temperature is still a meaningful coarse-grained variable, then for moderate temperature excursions we write
\begin{equation}
 du = \rho c_p \, dT,
\end{equation}
so that
\begin{equation}
 \frac{\partial u}{\partial t} \approx \rho c_p \frac{\partial T}{\partial t}.
\end{equation}

Thus energy conservation becomes
\begin{equation}
 \rho c_p \frac{\partial T}{\partial t} + \nabla\cdot \qvec = Q.
\end{equation}

\subsection{Phonon-aware constitutive law}

Now comes the difference from Module 4a. We still use a Fourier-like form for the heat flux,
\begin{equation}
 \qvec = -k(C,T)\nabla T,
\end{equation}
but the conductivity is no longer prescribed independently of damage. It is now a constitutive function that encodes phonon scattering by the evolving defect population.

Substituting this constitutive law into energy conservation gives
\begin{equation}
\rho c_p \frac{\partial T}{\partial t} = \nabla\cdot\left(k(C,T)\nabla T\right) + Q.
\end{equation}

This is the governing equation for Module 4b.

\subsection{Why this still differs from Module 4a even though the PDE looks similar}

At first sight, the equation resembles Module 4a. The difference is in the constitutive closure. In Module 4a, $k$ is a prescribed property or a simple known function, often independent of the evolving radiation damage field. In Module 4b, $k$ is updated by the damage state supplied by Module 3. This means Module 4b is not just a temperature solve; it is a defect-aware temperature solve.

This difference can change both the amplitude and spatial structure of the temperature field, especially if damage localizes in certain regions.

\section{Expansion of the conductivity-divergence term}

Because $k$ depends on both $C$ and $T$, the diffusion operator is no longer the simple constant-coefficient Laplacian. Expanding the divergence gives
\begin{equation}
\nabla\cdot\left(k\nabla T\right) = k\nabla^2 T + \nabla k \cdot \nabla T.
\end{equation}

This identity is important physically.

\begin{itemize}
    \item The term $k\nabla^2 T$ is the familiar diffusion-like piece.
    \item The term $\nabla k \cdot \nabla T$ appears because the conductivity varies in space.
\end{itemize}

Since $k=k(C,T)$, the conductivity gradient can be written by the chain rule as
\begin{equation}
\nabla k = \frac{\partial k}{\partial C}\nabla C + \frac{\partial k}{\partial T}\nabla T.
\end{equation}
Therefore,
\begin{equation}
\nabla\cdot\left(k\nabla T\right) = k\nabla^2 T + \frac{\partial k}{\partial C}\nabla C\cdot\nabla T + \frac{\partial k}{\partial T}\absval{\nabla T}^2.
\end{equation}

This expanded form shows why Module 4b is richer than Module 4a. Spatial nonuniformity in defect concentration directly enters the temperature equation through the $\partial k/\partial C$ term. Thus a localized defect cloud can alter not only the local magnitude of conductivity, but also the way temperature gradients propagate through space.

\section{Reduction to the two-dimensional project geometry}

The project implementation is two-dimensional in the post-Geant4 modules. Therefore, we restrict to coordinates $(x,y)$ and write
\begin{equation}
\rho c_p \frac{\partial T}{\partial t} = \frac{\partial}{\partial x}\left(k(C,T)\frac{\partial T}{\partial x}\right) + \frac{\partial}{\partial y}\left(k(C,T)\frac{\partial T}{\partial y}\right) + Q.
\end{equation}

If expanded explicitly,
\begin{align}
\rho c_p \frac{\partial T}{\partial t} = {} & k\left(\frac{\partial^2 T}{\partial x^2}+\frac{\partial^2 T}{\partial y^2}\right)
+ \frac{\partial k}{\partial x}\frac{\partial T}{\partial x}
+ \frac{\partial k}{\partial y}\frac{\partial T}{\partial y}
+ Q.
\end{align}

This is the form most directly tied to a 2D numerical implementation.

\section{Coupling of Module 4b to Module 3}

Module 3 evolves the defect field $C(\rvec,t)$ through diffusion, reaction, clustering, and annealing. Module 4b receives that field and uses it to update the conductivity. In return, the temperature field from Module 4b feeds back into Module 3 because defect kinetics are temperature dependent.

A common Arrhenius form for a temperature-dependent diffusion coefficient is
\begin{equation}
 D(T) = D_0 \exp\left(-\frac{E_a}{\kB T}\right),
\end{equation}
where $D_0$ is a prefactor and $E_a$ is an activation energy. Similarly, an annealing or reaction rate may follow
\begin{equation}
 k_{\mathrm{ann}}(T) = k_0 \exp\left(-\frac{E_{\mathrm{ann}}}{\kB T}\right).
\end{equation}

Therefore the full feedback loop is
\begin{equation}
 C \longrightarrow k(C,T) \longrightarrow T \longrightarrow D(T),\, k_{\mathrm{ann}}(T) \longrightarrow C.
\end{equation}

That feedback is one of the main scientific reasons this project becomes a multiphysics simulator rather than a loose collection of independent models.

\section{What Module 4b captures and what it does not capture}

\subsection{What it captures}

Module 4b captures the following physics beyond Module 4a:
\begin{enumerate}[label=\arabic*.]
    \item Thermal conductivity is not fixed independently of damage.
    \item Defect accumulation can locally suppress heat transport.
    \item Thermal transport becomes coupled to the evolving damage state.
    \item Spatial gradients in defect concentration can indirectly reshape temperature gradients.
\end{enumerate}

\subsection{What it does not yet capture}

Module 4b is still not a full microscopic phonon solver. In particular, in its first reduced form it does not explicitly resolve:
\begin{enumerate}[label=\arabic*.]
    \item spectral dependence of phonon transport,
    \item branch-resolved transport for longitudinal and transverse modes,
    \item strongly ballistic transport,
    \item full nonequilibrium phonon populations,
    \item detailed interface transmission physics.
\end{enumerate}

Thus Module 4b should be understood as a \textbf{phonon-informed continuum model}, not a full Boltzmann transport implementation.

\section{When Module 4b should outperform Module 4a}

Module 4b is likely to matter more than Module 4a when one or more of the following is true:
\begin{itemize}
    \item defect concentrations become large enough to noticeably suppress thermal conductivity,
    \item damage is spatially localized so that conductivity becomes strongly nonuniform,
    \item the device has thermal bottlenecks or hot spots whose severity depends sensitively on local conductivity,
    \item the project goal is to distinguish whether thermal feedback from damage materially changes long-term degradation predictions.
\end{itemize}

If none of these occur, Module 4a may already be adequate. Therefore the proper scientific use of Module 4b is not to assume it is automatically superior in every circumstance, but to compare it against Module 4a and determine whether the added physics changes the conclusions enough to justify the added complexity.

\section{Characteristic scales and nondimensional interpretation}

Suppose a characteristic device length scale is $L$, a characteristic temperature rise is $\Delta T$, and a characteristic conductivity scale is $k_0$. Define the nominal thermal diffusivity
\begin{equation}
 \alpha_{\mathrm{th},0} = \frac{k_0}{\rho c_p}.
\end{equation}
Then the nominal thermal diffusion time is
\begin{equation}
 t_{\mathrm{th},0} \sim \frac{L^2}{\alpha_{\mathrm{th},0}} = \frac{\rho c_p L^2}{k_0}.
\end{equation}

Now if defects reduce the conductivity to an effective value $k_{\mathrm{eff}}<k_0$, then the thermal diffusion time increases:
\begin{equation}
 t_{\mathrm{th,eff}} \sim \frac{\rho c_p L^2}{k_{\mathrm{eff}}} > t_{\mathrm{th},0}.
\end{equation}

This simple estimate already tells us the project-level consequence of damage-suppressed conductivity: thermal relaxation slows down, so hot regions can persist longer. Longer-lived hot regions can accelerate defect migration and annealing in a spatially selective way.

\section{Recommended first implementation strategy for Module 4b}

The most practical first implementation is not a BTE solver. The most practical first implementation is:
\begin{enumerate}[label=\arabic*.]
    \item Solve the same 2D thermal PDE as Module 4a.
    \item Replace constant or prescribed $k$ with a conductivity law $k(C,T)$.
    \item Start with one simple monotone reduction law, such as $k=k_{\mathrm{bulk}}(T)/(1+\alpha C)$.
    \item Compare the resulting temperature field against Module 4a for the same source and boundary conditions.
    \item Quantify differences in peak temperature, thermal relaxation time, and downstream defect evolution metrics.
\end{enumerate}

This strategy gives a clear and manageable path from baseline continuum transport to a more physically faithful defect-sensitive model.

\section{Verification and comparison logic}

The verification logic should proceed in layers.

\subsection{Internal consistency checks}
\begin{itemize}
    \item If $C\to 0$, Module 4b should reduce to the undamaged conductivity model.
    \item If $k(C,T)$ is artificially frozen to a constant, Module 4b should reduce to Module 4a.
    \item If the defect field is spatially uniform, then the conductivity field should also be spatially uniform, apart from temperature dependence.
\end{itemize}

\subsection{Comparative metrics against Module 4a}
Useful comparison outputs include:
\begin{itemize}
    \item peak temperature,
    \item spatial average temperature,
    \item time to cool below a specified threshold,
    \item temperature nonuniformity metric,
    \item defect-kinetics changes induced by the altered temperature history.
\end{itemize}

These are the natural quantities to show in preliminary slides when comparing the simpler and richer thermal models.

\section{Summary}

Module 4b is a phonon-aware thermal transport model for silicon that remains computationally compatible with the larger multiphysics project. Its central idea is that radiation-induced defects scatter phonons, and because phonons carry heat in silicon, the thermal conductivity should evolve with the defect population.

The derivation starts from the phonon Boltzmann transport equation only to establish the physical origin of conductivity. Under near-equilibrium, diffusive conditions, that microscopic theory reduces to a Fourier-like heat flux law. The key Module 4b step is then to replace the fixed conductivity of Module 4a with a defect- and temperature-dependent conductivity $k(C,T)$. This yields the governing equation
\begin{equation}
\rho c_p \frac{\partial T}{\partial t} = \nabla\cdot\left(k(C,T)\nabla T\right)+Q,
\end{equation}
which is the natural reduced phonon-informed extension of the Module 4a continuum model.

This model is richer than Module 4a because it embeds damage-induced thermal degradation into the transport law itself, yet it remains far cheaper and easier to couple than a full phonon BTE solver. That balance makes Module 4b the correct next-fidelity thermal model for the present project.

\end{document}
